\chapter{Het kanonieke ensemble}\label{ch:b3:canonical-ensemble}

Geïsoleerde stelsels zijn een verzinsel van theoretici: echte monsters
zitten in thermostaten, kamers en oceanen van lucht --- in contact met een
reservoir dat niet hun energie maar hun \emph{temperatuur} vastlegt. Het
volume van jaar 2 ontmoette de bijbehorende wet empirisch: de kans op een
toestand met energie $E$ draagt de factor $\eu^{-E/k_{\text{B}}T}$. Dit
hoofdstuk leidt die boltzmannfactor in vier regels af uit het tellen van
\cref{ch:b3:microcanonical-ensemble} en bouwt vervolgens de machine die
van de statistische fysica een industriële discipline maakt: de
\emph{\omterm{thm:b3:canonical-ensemble:boltzmann}{toestandssom}} $Z$, één som waaruit energie, entropie, druk,
warmtecapaciteit en fluctuaties alle door differentiëren vallen. De eerste
zeges van de machine worden hier naverteld: waarom warmtecapaciteiten bij
lage temperatuur sterven (Einstein, 1907 --- de eerste kwantumtheorie van
de materie), waarom de warmtecapaciteit van waterstof een trap opklimt,
waarom de hemel exponentieel ijler wordt --- en hoe Perrin, kijkend naar
microscopische korrels die in een druppel water bezonken, het getal van
Avogadro telde en het debat over het bestaan van atomen beëindigde.

\section{De boltzmannverdeling, afgeleid}

\begin{theorem}[Kanonieke verdeling]\label{thm:b3:canonical-ensemble:boltzmann}
\index{boltzmannverdeling}\index{kanoniek ensemble}\index{toestandssom}
Een stelsel dat energie uitwisselt met een groot reservoir op temperatuur
$T$ bezet zijn \omterm{def:b3:microcanonical-ensemble:states}{microtoestand} $s$ (energie $E_s$) met kans
\[
\mathcal P_s = \frac{\eu^{-E_s/k_{\text{B}}T}}{Z} , \qquad
Z = \sum_s \eu^{-E_s/k_{\text{B}}T}
\]
(voortaan $\beta = 1/k_{\text{B}}T$). $Z$, de \emph{\omterm{thm:b3:canonical-ensemble:boltzmann}{toestandssom}},
normeert de kansen --- en blijkt de volledige thermodynamica van het
stelsel te bevatten.
\end{theorem}

\begin{proof}
Het stelsel in toestand $s$ laat het reservoir de energie $E - E_s$: door
het \omterm{thm:b3:microcanonical-ensemble:postulate}{grondpostulaat} op het geheel toe te passen is $\mathcal P_s
\propto \Omega_{\text{res}}(E - E_s)$. Ontwikkel de entropie van het reservoir,
enorm en glad, tot eerste orde: $S_{\text{res}}(E - E_s)
= S_{\text{res}}(E) - E_s\,\partial S_{\text{res}}/\partial E =
S_{\text{res}}(E) - E_s/T$ volgens de definitie van temperatuur. Dus is
$\Omega_{\text{res}} \propto \eu^{-E_s/k_{\text{B}}T}$. (Hogere ordes
sterven met de omvang van het reservoir.)
\end{proof}

\begin{proposition}[De machine]\label{prop:b3:canonical-ensemble:machine}
\index{vrije energie}\index{warmtecapaciteit!en fluctuaties}
Uit $Z(T, V, N)$:
\[
\langle E\rangle = -\frac{\partial\ln Z}{\partial\beta} , \qquad
F = -k_{\text{B}}T\ln Z , \qquad
S = -\frac{\partial F}{\partial T} , \qquad
P = -\frac{\partial F}{\partial V} ,
\]
waarin $F = \langle E\rangle - TS$ de \emph{\omterm{prop:b3:canonical-ensemble:machine}{vrije energie}} is; en de
warmtecapaciteit is een \emph{fluctuatie}:
\[
C = \frac{\partial\langle E\rangle}{\partial T}
 = \frac{\langle E^2\rangle - \langle E\rangle^2}
{k_{\text{B}}T^2} .
\]
Voor onafhankelijke, onderscheidbare deelstelsels is $Z = z^N$; voor $N$
\omterm{def:b3:identical-particles:exchange}{identieke deeltjes} in één gemeenschappelijke doos $Z = z^N/N!$. Bij vaste
$T$ en $V$ \emph{minimaliseert} het evenwicht $F$: de natuur ruilt energie
tegen entropie tegen de wisselkoers $T$ --- het nuttigste beginsel uit de
hele fysica van de materie.
\end{proposition}

\begin{proof}[Gedeeltelijk bewijs]
$-\partial_\beta\ln Z = \sum E_s\eu^{-\beta E_s}/Z = \langle
E\rangle$;
nog een afgeleide geeft $\langle E^2\rangle - \langle
E\rangle^2$, en de
kettingregel zet $\partial_\beta$ om in $\partial_T$. De identificaties van
$F$, $S$ en $P$ volgen door $\dd(-k_{\text{B}}T\ln Z)$ met het
thermodynamische $\dd F = -S\dd T - P\dd V$ te vergelijken; het
factoriseren is de exponentiële functie van een som. Het minimumbeginsel:
een dalende $F_{\text{systeem}}$ is een stijgende $S_{\text{totaal}}$, want
$\Delta S_{\text{res}} =
-\Delta E_{\text{sys}}/T$.
\end{proof}

\begin{omfigure}
\begin{tikzpicture}[font=\small, >=Stealth]
  \fill[omExo!12, draw=omExo, thick, rounded corners=3pt] (-4.4,-1.9) rectangle (4.4,1.9);
  \node[omExo, anchor=north west, font=\footnotesize] at (-4.3,1.82)
    {reservoir op $T$: legt de wisselkoers vast};
  \fill[omDef!22, draw=omDef, very thick, rounded corners=2pt] (-1.05,-0.8) rectangle (1.05,0.8);
  \node[omDef] at (0,0.22) {stelsel};
  \node[omDef, font=\footnotesize] at (0,-0.3) {toestand $s$, $E_s$};
  \draw[<->, thick] (1.15,0) -- (2.5,0);
  \node[anchor=south, font=\footnotesize] at (1.85,0.08) {energie};
  \node[anchor=west, align=left, font=\footnotesize] at (5.0,0.6)
    {$\mathcal P_s \propto \eu^{-E_s/k_{\text{B}}T}$};
  \node[anchor=west, align=left, font=\footnotesize] at (5.0,-0.4)
    {kosten van $E_s$ lenen:\\reservoirentropie $-E_s/T$};
\end{tikzpicture}

{\small De kanonieke opstelling: een klein stelsel dat energie leent van
een reusachtig reservoir. Elke geleende joule kost het reservoir $1/T$ aan
entropie --- vandaar de exponentiële korting op energierijke toestanden.}
\end{omfigure}

\section{Eerste zeges}

\begin{example}[Twee niveaus, in twee regels]\label{ex:b3:canonical-ensemble:twolevel}
Voor de niveaus $0, \epsilon$ is $z = 1 + \eu^{-\beta\epsilon}$, wat
$\langle E\rangle = N\epsilon/(\eu^{\beta\epsilon} + 1)$ geeft --- het
resultaat dat de microkanonieke weg een bladzijde Stirling kostte
(\cref{exo:b3:microcanonical-ensemble:8}). Het kanonieke formalisme is het
microkanonieke met de combinatoriek voorgekauwd; de schottkybult in $C(T)$
volgt uit één differentiatie.
\end{example}

\begin{example}[De vaste stof van Einstein en de dood van Dulong--Petit]\label{ex:b3:canonical-ensemble:einstein}
\index{vaste stof van Einstein}\index{wet van Dulong--Petit}
Modelleer een kristal als $3N$ kwantumoscillatoren met frequentie $\omega$
(\cref{ch:b3:harmonic-oscillator}). Per oscillator geldt
\[
z = \sum_n\eu^{-\beta\hbar\omega(n + 1/2)}
 = \frac{1}{2\sinh(\beta\hbar\omega/2)} , \qquad
\langle E\rangle = \frac{\hbar\omega}{2} +
\frac{\hbar\omega}{\eu^{\beta\hbar\omega} - 1} .
\]
Bij hoge $T$ gaat $\langle E\rangle \to k_{\text{B}}T$ per oscillator en
$C \to 3Nk_{\text{B}}$ --- de wet van Dulong en Petit uit het volume van
jaar 1, verklaard. Bij lage $T$ wordt het kwantum $\hbar\omega$
onbetaalbaar en stort $C$ exponentieel in --- zoals gemeten, en klassiek
onverklaarbaar. De kromme van Einstein uit 1907, met één parameter per
element, was de eerste toepassing van quanta op gewone materie; diamant,
met stijve bindingen en lichte atomen ($\hbar\omega/k_{\text{B}}
\approx \qty{1300}{K}$), is bij kamertemperatuur nog steeds ``bevroren'' --- de
anomalie die scheikundigen tachtig jaar had geplaagd
(\cref{exo:b3:canonical-ensemble:6}).
\end{example}

\begin{omfigure}
\begin{tikzpicture}
  \begin{axis}[omaxis, axis x line=bottom, axis y line=left,
    width=6.8cm, height=5.2cm, xmin=0, xmax=2, ymin=0, ymax=1.15,
    xlabel={$T/\theta_{\text{E}}$}, ylabel={$C/3Nk_{\text{B}}$},
    xlabel style={at={(axis description cs:0.5,-0.11)}, anchor=north},
    xtick={1,2}, ytick={1}, clip=false]
    \addplot[gray, dashed, domain=0:2] {1};
    \addplot[omDef, very thick, domain=0.02:2, samples=200]
      {(1/x)^2*exp(1/x)/(exp(1/x)-1)^2};
    \node[font=\scriptsize, anchor=west, gray] at (axis cs:0.08,1.07)
      {Dulong--Petit};
    \node[omDef, font=\scriptsize, anchor=west] at (axis cs:0.75,0.5)
      {Einstein 1907};
  \end{axis}
  \begin{axis}[omaxis, axis x line=bottom, axis y line=left,
    width=6.8cm, height=5.2cm, xmin=1, xmax=4000, ymin=0, ymax=4.3,
    xshift=7.4cm, xmode=log,
    xlabel={$T$ (K, logschaal)}, ylabel={$C_V/Nk_{\text{B}}$ voor H$_2$},
    xlabel style={at={(axis description cs:0.5,-0.11)}, anchor=north},
    xtick={10,100,1000}, xticklabels={10,100,1000},
    ytick={1.5,2.5,3.5}, clip=false]
    \addplot[omThm, very thick, domain=2:4000, samples=400]
      {1.5 + (87/x)^2*exp(-87/x)/(1-exp(-87/x))^2
       + ((6000/x)<25 ? (6000/x)^2*exp(-6000/x)/(1-exp(-6000/x))^2 : 0)};
    \node[omThm, font=\scriptsize, anchor=west] at (axis cs:3,1.68) {translatie};
    \node[omThm, font=\scriptsize, anchor=west] at (axis cs:150,2.72) {+ rotatie};
    \node[omThm, font=\scriptsize, anchor=south east] at (axis cs:3800,3.55) {+ trilling};
  \end{axis}
\end{tikzpicture}

{\small Links: de warmtecapaciteitskromme van Einstein --- de klassieke
equipartitie teruggewonnen bij hoge $T$, kwantumbevriezing onder
$\theta_{\text{E}} = \hbar\omega/k_{\text{B}}$. Rechts: de $C_V$ van
waterstofgas klimt een trap op naarmate de rotatie (nabij \qty{85}{K}) en
daarna de trilling (nabij \qty{6000}{K}, buiten beeld) ontdooien: elke
beweging doet pas aan de equipartitie mee wanneer $k_{\text{B}}T$ haar
kwantum kan betalen.}
\end{omfigure}

\begin{theorem}[Equipartitie, met haar vergunning]\label{thm:b3:canonical-ensemble:equipartition}
\index{equipartitiestelling}
Elke coördinaat of impuls die \emph{kwadratisch} in de energie voorkomt
draagt in het klassieke regime (bij hoge temperatuur)
\[
\langle\epsilon\rangle = \tfrac12 k_{\text{B}}T
\]
aan de gemiddelde energie bij: $\tfrac32 k_{\text{B}}T$ voor een atoom van
een eenatomig gas, $\tfrac52$ voor een roterend tweeatomig molecuul en
$3k_{\text{B}}T$ voor een oscillator. De vergunning verloopt zodra
$k_{\text{B}}T$ onder de niveauafstand van de modus zakt: de modus vriest
uit en haar bijdrage verdwijnt --- de oplossing van de
warmtecapaciteitsschandalen van de negentiende eeuw, hierboven als trap
getekend.
\end{theorem}

\begin{proof}[Gedeeltelijk bewijs]
Voor $\epsilon = ax^2$ geldt
\[
\langle\epsilon\rangle
 = \frac{\int ax^2\,\eu^{-\beta ax^2}\dd x}
        {\int\eu^{-\beta ax^2}\dd x}
 = -\partial_\beta\ln\!\int\eu^{-\beta ax^2}\dd x
 = -\partial_\beta\ln\beta^{-1/2} = \frac{1}{2\beta} .
\] Het bevriezen is de berekening van
\cref{ex:b3:canonical-ensemble:einstein}, modus voor modus.
\end{proof}

\begin{example}[Het gas, kanoniek]\label{ex:b3:canonical-ensemble:gas}
Eén atoom in een doos: $z = V/\lambda_T^3$ (de toestandstelling van
\cref{prop:b3:schrodinger-three-dimensions:counting}, met Boltzmann
gewogen); $N$ identieke atomen: $Z = z^N/N!$. Dan is $F =
-Nk_{\text{B}}T[\ln(V/N\lambda_T^3) + 1]$, en differentiëren levert $PV = Nk_{\text{B}}T$, de entropie van Sackur en Tetrode, en $\langle E\rangle = \tfrac32 Nk_{\text{B}}T$ --- het hele ideale gas uit één gaussische
integraal. De boltzmannfactor toegepast op de kinetische energie van een
molecuul geeft de snelheidsverdeling van Maxwell uit het volume van jaar
1, nu afgeleid; toegepast op zijn potentiële energie $mgh$ geeft hij de
exponentiële atmosfeer --- en, in een druppel water, de korrelladder van
Perrin (\cref{pb:b3:canonical-ensemble:1}).
\end{example}

\begin{method}[Kanoniek ambacht]\label{met:b3:canonical-ensemble:method}
(1) Noem de toestanden en energieën van \emph{één} eenheid; bereken $z$.
(2) Factoriseer: $Z = z^N$ (of $z^N/N!$ voor \omterm{def:b3:identical-particles:exchange}{identieke deeltjes} die de
ruimte delen); neem vroeg de $\ln$. (3) Differentieer: naar $\beta$ voor de
energie, naar $T$ voor de entropie via $F$, naar $V$ voor de druk; een
tweede afgeleide voor $C$ en de fluctuaties. (4) Toets beide uiteinden:
hoge $T$ moet de equipartitie opleveren, lage $T$ moet met een staart
$\eu^{-\Delta/k_{\text{B}}T}$ bevriezen. (5) Wedstrijden (vouwen, binden,
uitlijnen): schrijf $F = E - TS$ voor elk alternatief op en laat de
kleinste winnen --- het omslagpunt ligt bij $T \approx \Delta
E/\Delta S$.
\end{method}

\section{Opgaven}

\begin{exercise}[$\star$]\label{exo:b3:canonical-ensemble:1}
(a) Schrijf de bezettingsverhouding op van twee niveaus die $\epsilon$
uiteen liggen bij temperatuur $T$. (b) Welk deel van de natriumatomen zit
in een vlam van \qty{2500}{K} in de aangeslagen toestand van \qty{2.1}{eV}
van de D-lijn (ontaarding 2 voor de grondtoestand, 6 voor de aangeslagen:
bezettingsverhouding $3\eu^{-\beta\epsilon}$)?
(c) Waarom straalt de vlam desondanks fel geel (hoeveel atomen per
cm$^3$ volstaan)? (d) Bij welke temperatuur zou de aangeslagen fractie
$10\%$ halen?
\end{exercise}

\begin{exercise}[$\star$]\label{exo:b3:canonical-ensemble:2}
Een stelsel met drie niveaus: $0, \epsilon, 2\epsilon$. (a) Schrijf $z$
op. (b) Bereken $\langle E\rangle$ en toets beide temperatuurlimieten. (c)
Bij welke $T$ is het middelste niveau \emph{in absolute zin} het sterkst
bezet? (d) Toon aan dat geen enkele temperatuur, hoe hoog ook, een hoger
niveau sterker bezet maakt dan een lager --- welk begrip uit hoofdstuk 16
zou daarvoor nodig zijn?
\end{exercise}

\begin{exercise}[$\star$]\label{exo:b3:canonical-ensemble:3}
De isotherme atmosfeer. (a) Pas de boltzmannfactor toe op de potentiële
energie $mgh$ en leid $n(h) = n_0\eu^{-mgh/k_{\text{B}}
T}$ af. (b) Bereken
de schaalhoogte voor lucht ($m = \qty{4.8e-26}{kg}$) bij \qty{288}{K}. (c)
De druk op de top van de Mount Everest als fractie van die op zeeniveau.
(d) Waarom is de afname van de echte atmosfeer bijna maar niet precies
exponentieel (wat hebben we constant gehouden dat het niet is)?
\end{exercise}

\begin{exercise}[$\star$]\label{exo:b3:canonical-ensemble:4}
Boekhouding met equipartitie. Tel de kwadratische termen en voorspel de
molaire $C_V$ van (a) argon; (b) N$_2$ bij kamertemperatuur (rotatie aan,
trilling bevroren); (c) N$_2$ bij \qty{3000}{K}; (d) een kristallijne
vaste stof (Dulong--Petit). Waar komen de gemeten waarden $12.5$, $20.8$,
$\approx26$ en $\approx\qty{25}{J/(mol.K)}$ overeen, en wat leert elke
afwijking?
\end{exercise}

\begin{exercise}[$\star\star$]\label{exo:b3:canonical-ensemble:5}
De kwantumoscillator, kanoniek. (a) Sommeer de meetkundige reeks voor $z$.
(b) Leid $\langle E\rangle$ af en herken $\langle
n\rangle = 1/(\eu^{\beta\hbar\omega} - 1)$ --- het resultaat van
\cref{exo:b3:harmonic-oscillator:6}, nu moeiteloos. (c) Differentieer voor
$C(T)$ en ga de twee limieten na. (d) Waarom staat de vorm van $\langle n\rangle$ op het punt beroemd te worden
(\cref{ch:b3:photons-phonons})?
\end{exercise}

\begin{exercise}[$\star\star$]\label{exo:b3:canonical-ensemble:6}
Einstein tegen de gegevens. (a) Bij welke $T/\theta_{\text{E}}$ is $C$
volgens de formule van de figuur tot de helft van Dulong--Petit gezakt?
(b) Diamant: $\theta_{\text{E}} \approx \qty{1300}{K}$ --- bereken
$C/3Nk_{\text{B}}$ bij \qty{300}{K} en verklaar de ``anomalie van diamant''
uit de negentiende eeuw. (c) Lood: $\theta_{\text{E}}
\approx \qty{90}{K}$
--- waarom gedroeg lood zich altijd ``keurig''? (d) Metingen bij zeer lage
$T$ tonen $C \propto T^3$ en geen exponentiële afname: welke aanname van
Einstein faalt (alle oscillatoren \emph{één} frequentie), en wie
repareerde haar (\cref{ch:b3:photons-phonons})?
\end{exercise}

\begin{exercise}[$\star\star$]\label{exo:b3:canonical-ensemble:7}
De staart van Maxwell en de ontbrekende waterstof. (a) Schrijf uit de
boltzmannfactor op de kinetische energie de snelheidsverdeling op en
bepaal de waarschijnlijkste snelheid voor N$_2$ en H$_2$ bij \qty{288}{K}.
(b) De ontsnappingssnelheid van de aarde is \qty{11.2}{km/s}: bereken
$mv_{\text{esc}}^2/2k_{\text{B}}T$ voor beide gassen. (c) De ontsnappende
fractie gaat als $\eu^{-mv_{\text{esc}}^2/2k_{\text{B}}T}$: vergelijk de
twee exponenten en concludeer welk gas op geologische tijdschaal weglekt.
(d) Breng dit in verband met de waargenomen samenstelling van de atmosfeer
van de aarde (geen vrije H$_2$) tegenover die van Jupiter (grotendeels
H$_2$) --- welke twee parameters draaien het oordeel om?
\end{exercise}

\begin{exercise}[$\star\star$]\label{exo:b3:canonical-ensemble:8}
Fluctuaties ontmoeten de respons. (a) Bewijs $C = (\langle E^2\rangle -
\langle E\rangle^2)/k_{\text{B}}T^2$ uit twee afgeleiden van $\ln
Z$. (b)
Ga haar expliciet na op het stelsel met twee niveaus. (c) Toon voor $N$
onafhankelijke eenheden aan dat de \emph{relatieve} energiefluctuatie als
$1/\sqrt N$ daalt. (d) Formuleer de moraal: wat heeft de bereidheid van een
stelsel om warmte op te nemen (een respons) te maken met hoezeer zijn
energie trilt (een fluctuatie) --- de terugkerende ruilhandel van de
statistische fysica.
\end{exercise}

\begin{exercise}[$\star\star$]\label{exo:b3:canonical-ensemble:9}
Boltzmann in het scheikundelab. Reactietempo's dragen de factor
$\eu^{-E_a/k_{\text{B}}T}$ (het overwinnen van een activeringsbarrière
$E_a$ --- Arrhenius). (a) Toon aan dat de vuistregel ``het tempo
verdubbelt om de \qty{10}{K} nabij kamertemperatuur'' overeenkomt met $E_a \approx
\qty{0.55}{eV}$. (b) Met welke factor vertraagt die reactie in een
koelkast (\qty{5}{\celsius})? (c) Een ei koken op hoogte: op een pas van
\qty{2000}{m} kookt water bij \qty{93}{\celsius} --- schat de extra
kooktijd. (d) Waarom regeert de staart voorbij de barrière de scheikunde
en niet de gemiddelde energie (welke moleculen reageren)?
\end{exercise}

\begin{exercise}[$\star\star\star$]\label{exo:b3:canonical-ensemble:10}
Vouwen met twee toestanden. Een biomolecuul is gevouwen (energie $0$, één
configuratie) of ontvouwen (energie $\Delta E > 0$, $\Omega_u =
\eu^{\Delta S/k_{\text{B}}}$ configuraties). (a) Schrijf de gevouwen
fractie tegen $T$ op. (b) Toon aan dat het ``smelt''midden bij
$T_{\text{m}} = \Delta E/\Delta S$ ligt en duid het als een gelijkspel $F = E -
TS$. (c) Bereken met $\Delta E = \qty{3.0}{eV}$ en $\Delta S =
100\,k_{\text{B}}$ (de coöperatieve eenheid van een klein eiwit) de waarde
van $T_{\text{m}}$ en de breedte van de overgang. (d) Waarom verscherpt
\emph{coöperativiteit} (vele contacten die samen breken, een grote $\Delta E$ \emph{en} $\Delta S$) het smelten --- en hoe buiten machines voor
thermische cycli van DNA (PCR) precies dit uit?
\end{exercise}

\begin{exercise}[$\star\star\star$]\label{exo:b3:canonical-ensemble:11}
Paramagnetisme en magnetisch koelen. $N$ spins $\tfrac12$ met moment $\mu$
in een veld $B$. (a) Toon aan dat $M = N\mu\tanh(\mu B/k_{\text{B}}T)$ en
ontwikkel tot de wet van Curie $M \approx N\mu^2B/k_{\text{B}}T$. (b)
Reken de uitlijning $\mu B/k_{\text{B}}T$ uit voor elektronmomenten bij $B = \qty{1}{T}$ en $T = \qty{300}{K}$, en bij \qty{1}{K}. (c) Adiabatische
demagnetisatie: magnetiseer bij \qty{1}{K}, isoleer, en verlaag $B$
tienvoudig --- toon aan dat een constante entropie een constante $\mu B/k_{\text{B}}T$ betekent, zodat $T$ tienvoudig zakt. (d) Wat legt de
bodem van deze koelkast vast (de wisselwerkingen tussen de spins --- schat
de dipolaire schaal $\mu_0\mu^2/4\pi a^3$ voor $a =
\qty{0.5}{nm}$, in
temperatuureenheden)?
\end{exercise}

\begin{exercise}[$\star\star\star$]\label{exo:b3:canonical-ensemble:12}
De rotatietrap, kwantitatief. De rotatieniveaus van een tweeatomig
molecuul zijn $E_J = BJ(J+1)$, met ontaarding $2J + 1$
(\cref{ch:b3:quantum-angular-momentum}). (a) Schrijf $z_{\text{rot}}$ op
en toon aan dat de som voor $k_{\text{B}}T \gg B$ de integraal
$k_{\text{B}}T/B$ wordt: de $k_{\text{B}}$ van de equipartitie in $C_V$
teruggewonnen. (b) Definieer $\theta_{\text{rot}} =
B/k_{\text{B}}$ en
reken haar uit voor H$_2$ ($B = \qty{7.5}{meV}$) en N$_2$ ($B = \qty{0.25}{meV}$): welk gas vertoont de rotatiestap bij toegankelijke
temperaturen? (c) Schets de volledige trap $C_V(T)$ van waterstof met haar
drie plateaus en twee stijgingen, met getallen bij beide. (d) Het gemeten
gedrag bij lage $T$ van H$_2$ wordt nog gecompliceerd door het mengsel
ortho--para van 3:1 uit \cref{exo:b3:identical-particles:8}: zeg in één zin
hoe de statistiek van de kernspins tot in de warmtecapaciteit van een gas
reikt.
\end{exercise}

\section{Vraagstuk: Avogadro tellen in een druppel water}

\begin{problem}[{Weekendvraagstuk --- de korrels van Perrin en het
bestaan van atomen}]\label{pb:b3:canonical-ensemble:1}
In 1908 bracht Jean Perrin microscopische harskorrels in water in
zwevende toestand, liet ze bezinken en telde ze onder de microscoop laag
voor laag. De exponentiële ladder die hij vond was de boltzmannfactor
zichtbaar gemaakt --- en uit haar schaalhoogte haalde hij het getal van
Avogadro, waarmee hij de laatste sceptici overtuigde dat atomen bestaan.
Nobelprijs, 1926. Gegevens: gutti-korrels met straal $r =
\qty{0.212}{\micro m}$ en dichtheid $\rho_{\text{g}} =
\qty{1207}{kg/m^3}$; water $\rho_{\text{w}} = \qty{999}{kg/m^3}$;
$T = \qty{293}{K}$; $g = \qty{9.81}{m/s^2}$.

\medskip\noindent\textbf{Deel I --- De zichtbare boltzmannfactor.}
\begin{enumerate}
  \item Een korrel in water voelt de zwaartekracht min de opwaartse
        kracht: bereken haar \emph{effectieve} massa $m' = \tfrac43\pi r^3(\rho_{\text{g}}
        - \rho_{\text{w}})$ en haar gewicht.
  \item Schrijf het evenwichtsprofiel van de concentratie $n(h)$ op uit de
        boltzmannfactor.
  \item Bereken de schaalhoogte $h_0 = k_{\text{B}}T/m'g$ met de moderne
        $k_{\text{B}}$.
  \item Waarom moeten de korrels zo nauwkeurig even groot zijn (hoe hangt
        $h_0$ van $r$ af)?
  \item Vergelijk $h_0$ met dezelfde formule voor luchtmoleculen: waarom
        is de atmosfeer van de korrels micrometers hoog en die van de lucht
        kilometers?
  \item Perrin telde (in één reeks) de relatieve concentraties
        $100 : 55 : 30 : 17$ op vier gelijk verdeelde diepten met
        \qty{30}{\micro m} ertussen: ga na dat dit een exponentiële ladder
        is en haal haar $h_0$ eruit.
\end{enumerate}

\medskip\noindent\textbf{Deel II --- Het onzichtbare wegen.}
\begin{enumerate}[resume]
  \item Keer om: haal uit de gemeten $h_0$ uit punt 6 en de bekende $m'g$
        de waarde $k_{\text{B}}$.
  \item De gasconstante $R = \qty{8.314}{J/(mol.K)}$ was uit de
        macroscopische scheikunde bekend: combineer haar met uw
        $k_{\text{B}}$ om het getal van Avogadro $N_{\text{A}} =
        R/k_{\text{B}}$ te krijgen.
  \item De reeksen van Perrin gaven $N_{\text{A}}$ tussen $5.5$ en
        $\num{7.2e23}$: vergelijk met de moderne \num{6.022e23} en geef
        commentaar op die prestatie, gezien zijn microscoop en stopwatch.
  \item Bereken uit $N_{\text{A}}$ de massa van één \omterm{thm:b3:hydrogen-atom:levels}{waterstofatoom} --- het
        getal waar de atomisten een eeuw naar hadden verlangd.
  \item Leg de logische opbouw uit: welke \emph{macroscopische} metingen
        ($R$; de korrelgrootte en -dichtheid; de ladder) worden gecombineerd
        om één atoom te wegen, zonder dat er ooit een atoom is gezien?
  \item Waarom gehoorzamen de korrels --- elk $10^{10}$ atoommassa's ---
        aan dezelfde boltzmannstatistiek als moleculen (wat in de afleiding
        van \cref{thm:b3:canonical-ensemble:boltzmann} bekommert zich om
        de afmeting)?
\end{enumerate}

\medskip\noindent\textbf{Deel III --- De trilling die het bezegelt.}
\begin{enumerate}[resume]
  \item Dezelfde korrels trillen: de formule van Einstein uit 1905 voor de
        brownse beweging geeft $\langle x^2\rangle = 2Dt$ met $D =
        k_{\text{B}}T/6\pi\eta r$ (water: $\eta =
        \qty{e-3}{Pa.s}$).
        Bereken $D$ voor de korrels van Perrin.
  \item Hoe ver dwaalt een korrel in één minuut af? Kon Perrin dat met een
        micrometeroculair en een stopwatch meten?
  \item Perrin ging $\langle x^2\rangle \propto t$ na en haalde er
        \emph{opnieuw} en onafhankelijk $k_{\text{B}}$ uit: waarom droegen
        twee ongerelateerde wegen (een statische ladder; een dynamische
        trilling) naar één getal zoveel bewijskracht?
  \item De trilling is de equipartitie toegepast op de korrel: elke
        snelheidscomponent draagt $\tfrac12 k_{\text{B}}T$. Schat de
        effectieve thermische snelheid van de korrel ($m \approx
        \qty{4.8e-17}{kg}$).
  \item Waarom wordt die snelheid nooit rechtstreeks gezien (wat
        onderbreekt de vrije vlucht al na nanometers), en wat ziet men
        \emph{wel}?
  \item Zeg welke twee hoofdstukken van dit boek elkaar in de waarneming
        ontmoeten: de mechanica van de wrijving (volume van jaar 2,
        viscositeit) en de statistiek van dit hoofdstuk.
\end{enumerate}

\medskip\noindent\textbf{Deel IV --- Wat er was beslecht.}
\begin{enumerate}[resume]
  \item Ostwald en Mach hielden atomen voor boekhoudkundige verzinsels:
        zeg in één zin waarom een \emph{getelde} $N_{\text{A}}$ uit
        korrelladders die stelling beëindigde.
  \item Noem drie andere wegen uit de jaren 1900 die op dezelfde
        $N_{\text{A}}$ uitkwamen (het blauw van de hemel,
        \cref{exo:b3:microcanonical-ensemble:12}; elektrolyse plus de
        elektronlading; de heliumproductie van radioactiviteit) --- waarom
        telde de \emph{convergentie} zwaarder dan elke afzonderlijke
        waarde?
  \item De ladder van Perrin is een evenwicht tussen welke twee munten van
        dit hoofdstuk (energie die omlaag trekt, entropie die uitspreidt),
        beprijsd tegen welke koers?
  \item Moderne toepassingen van dezelfde natuurkunde: analytische
        ultracentrifuges laten eiwitten bij $10^5g$ ronddraaien om hun
        ``atmosferen'' tot meetbare ladders samen te persen --- toon aan dat
        $g$ met $10^5$ vermenigvuldigen $h_0$ met dezelfde factor deelt, en
        schat $h_0$ voor een eiwit met effectieve massa \qty{e-22}{kg} bij
        $10^5g$ en \qty{293}{K}.
  \item De korrelstraal verdubbelen deelt $h_0$ door acht: begrens het
        praktische venster van korrelgrootten tussen ``ladder te hoog om
        een gradiënt te zien'' en ``ladder dunner dan één korrel'' voor een
        microscoopveld van \qty{100}{\micro m} diep.
  \item Sinds 2019 \emph{definieert} het SI $k_{\text{B}}$ en
        $N_{\text{A}}$ exact: zeg wat een proef in de stijl van Perrin
        \emph{vandaag} meet (een consistentietoets, of een ijking van het
        toestel en de korrels) --- en waarom de natuurkunde ongewijzigd is.
  \item Vat het genoemde resultaat samen: de concentratie van een korrel
        van \qty{0.2}{\micro m} halveert om de $\sim\qty{35}{\micro m}$
        hoogte; gelezen met $n(h) = n_0\eu^{-m'gh/k_{\text{B}}T}$ leverde
        die ladder $N_{\text{A}} \approx \num{6e23}$ --- atomen geteld en
        niet vermoed, in een druppel water op een microscooptafel.
\end{enumerate}
\end{problem}
